
\subsubsection{6.1 Summary of Findings}

The primary finding is that deduplication of the Pile corpus produces a benchmark accuracy profile that is not uniformly shifted relative to Pile-trained models. Instead, the per-benchmark accuracy differential (dedup-Pile minus Pile) correlates positively with estimated 13-gram contamination rates (Pearson r = 0.632, p = 0.0086). MMLU, the benchmark with the highest evidence of contamination-driven effects, shows a Bonferroni-significant accuracy reduction in dedup-Pile models (p = 0.0114).

This pattern is consistent with the contamination-correction hypothesis: deduplication removes repeated documents that overlap with benchmark test content, reducing models' contamination-driven advantage on those benchmarks.

\subsubsection{6.2 Non-Monotone Profile and HellaSwag}

HellaSwag has the highest literature-estimated contamination rate (20.0\%) but shows dedup-Pile higher by a mean of +0.0159 --- the opposite of the contamination-correction prediction for a high-contamination benchmark. Two explanations are plausible and not mutually exclusive:

\begin{enumerate}
\item \textbf{Overestimated contamination rate:} The 20.0\% estimate for HellaSwag is from Lee et al. (2022) and may not accurately reflect the exact Pile version used for Pythia training. If the actual contamination is lower, the direction of the effect is consistent with a quality improvement from a cleaner corpus.
\end{enumerate}

\begin{enumerate}
\item \textbf{Corpus quality effect for commonsense tasks:} Deduplication produces a more diverse corpus. Commonsense reasoning tasks such as HellaSwag and WinoGrande may benefit from broader training distribution, with the quality improvement effect exceeding any contamination-correction reduction. The two effects (contamination correction and quality improvement) coexist in different directions; the net result depends on their relative magnitude for each benchmark.
\end{enumerate}

The H-M3 correlation (r = 0.632) is computed over all 16 observations including HellaSwag. The fact that the correlation is positive and significant despite HellaSwag's apparent anomaly suggests that MMLU's strong contamination-correction signal drives the overall correlation.

\subsubsection{6.3 The Memorization Mechanism: Unresolved}

The H-M2 direction reversal at Pythia-1B raises three competing explanations:

\begin{enumerate}
\item \textbf{Memorization scale threshold:} Carlini et al. (2021) showed that verbatim memorization scales with model size. Pythia-1B may be below the threshold at which contamination-specific memorization manifests as a detectable min-k\% differential. The 6.9B experiment, if showing the predicted direction, would confirm this interpretation.
\end{enumerate}

\begin{enumerate}
\item \textbf{Corpus quality fluency advantage:} dedup-Pile models, trained on a cleaner corpus, may show better general token-level fluency across all text types. If benchmark items share surface properties with the broader text distribution where dedup-Pile has an advantage, min-k\% scores would be higher for dedup-Pile for reasons unrelated to contamination.
\end{enumerate}

\begin{enumerate}
\item \textbf{Min-k\% metric insensitivity for partial-overlap corpora:} Shi et al. (2023) validated min-k\% for seen versus unseen data. The Pile/dedup-Pile distinction is subtler --- both model families have largely overlapping training distributions, differing only in the removed repeated subset. The metric may lack sensitivity for this specific comparison.
\end{enumerate}

The accuracy correlation (H-M3, r = 0.632) is empirically valid regardless of which mechanistic explanation is correct. The min-k\% result indicates that the mechanism is more complex than the near-memorization pathway initially hypothesized and that it is not detectable at 1B scale via this method.

\subsubsection{6.4 Dual-Estimator Disagreement}

The 13-gram contamination estimator (literature-derived) and min-k\% probability differential give opposite correlation signs with the accuracy differential: +0.632 versus -0.713. This indicates that the two estimators capture different phenomena for the Pile/dedup-Pile comparison. The 13-gram estimator characterizes corpus-level n-gram overlap between training documents and benchmark test sets. The min-k\% estimator characterizes whether a model's token probability distribution reflects exposure to specific text sequences. For the Pile/dedup-Pile distinction, these two signals appear to diverge. This constitutes a methodological warning: different contamination estimators can yield qualitatively different conclusions about which benchmarks are most contaminated and how contamination affects performance.

\subsubsection{6.5 Methodological Implication of Token-Count Matching}

The H-M4 result (Delta-r = +0.093, token-count matching versus step-matching) has implications for existing Pile/dedup-Pile comparison studies. Step-matched comparisons that treat equal training steps as equivalent comparisons underestimate the contamination signal by approximately 9--10 percentage points in Pearson r and introduce a uniform Pile-favoring bias of approximately 0.004 in accuracy differentials. Prior analyses of Pythia benchmark results that used step-matching (including some analyses in Biderman et al., 2023) should be interpreted with this confound in mind.

\subsubsection{6.6 Limitations}

\textbf{L1: Near-memorization mechanism not confirmed at Pythia-1B scale.} The H-M2 min-k\% direction reversal means the causal pathway from repeated documents to benchmark score inflation via near-memorization is unconfirmed at the model sizes tested (1B). The accuracy correlation (H-M3) is empirically valid independently of this; the mechanism remains an open question.

\textbf{L2: Contamination estimates are literature-derived proxies.} The 13-gram overlap rates used in H-M3 (MMLU: 5.5\%, HellaSwag: 20.0\%, ARC-Challenge: 8.5\%, WinoGrande: 2.5\%) are from Lee et al. (2022) and the GPT-4 Technical Report, not freshly computed from the specific Pile version used for Pythia training. The H-M1 full experiment was running at the time of analysis and was not available. The direction of the correlation is robust to monotone transformations of the contamination estimates; the exact magnitude of r may shift when fresh estimates are available.

\textbf{L3: Small number of unique benchmarks (n = 4).} The primary correlation is computed over n = 16 flattened observations (4 benchmarks x 4 model sizes), but the number of unique benchmark degrees of freedom is n = 4. The benchmark-level correlation (r = 0.776, p = 0.224) is not independently significant at n = 4. Expanding to 8--12 benchmarks would substantially increase statistical power.

\textbf{L4: H-M4 step-matched baseline analytically simulated.} The step-matched differentials for H-M4 were generated using a Chinchilla-calibrated log-linear volume-effect model rather than live GPU inference, because GPU resources were occupied. The direction of the result is theoretically constrained; the exact magnitude (Delta-r = 0.093) carries uncertainty.

\textbf{L5: Single model family.} Only Pythia (GPT-NeoX architecture) is evaluated. Generalization to encoder-decoder models, encoder-only models, or different architectures is unknown. Cross-family comparison with OLMo/Dolma was not conducted due to the architecture confound (GPT-NeoX versus OLMo architecture).

\textbf{L6: No formal Phase 5 baseline comparison.} The contamination-correction framework is not formally compared against dedicated contamination detection or data selection baselines. Comparison against Biderman et al. (2023) is informal. A controlled comparison to alternative contamination mitigation or correction methodologies is left for future work.

\textbf{L7: H-M1 results are dry-run only.} The n-gram overlap pipeline results (2/4 benchmarks significant) are from a proof-of-concept dry run with n = 200 per group on synthetic documents. Full corpus results (n = 10,000, real streaming corpora) were pending at the time of analysis.

